\documentclass[12pt,a4paper]{article}

\usepackage[T2A]{fontenc}
\usepackage[utf8]{inputenc}
\usepackage[russian]{babel}
\usepackage{amsmath,amssymb}
\usepackage{array}
\usepackage{longtable}
\usepackage{enumitem}
\usepackage{geometry}
\usepackage{hyperref}

\geometry{left=25mm,right=20mm,top=20mm,bottom=20mm}

\title{Билет №57 \\ \small Терминология: «программная инженерия» и «технология программирования». Причины возникновения программной инженерии как науки. (СпБГУ) \\ Что такое программное изделие и программный продукт. Основные отличия промышленного, Open source и исследовательского программного приложения. (СпБПУ)}
\author{Сериков Владислав Александрович}
\date{16 мая 2026}

\begin{document}

\maketitle
\tableofcontents
\newpage

\section{Программная инженерия и технология программирования}

\subsection{Понятие программной инженерии}

\textbf{Программная инженерия} --- это область инженерной деятельности и научно-практическая дисциплина, изучающая методы, процессы, средства и организационные принципы создания, эксплуатации, сопровождения и развития программного обеспечения.

Классическое определение, принятое в стандартах IEEE, формулирует программную инженерию как применение систематического, дисциплинированного и количественно измеримого подхода к разработке, эксплуатации и сопровождению программного обеспечения, то есть применение инженерии к программному обеспечению.

Из этого определения следуют три важных признака программной инженерии:

\begin{enumerate}
    \item \textbf{Систематичность}: разработка программ не рассматривается как набор случайных действий программиста; она организуется как процесс с этапами, правилами, ролями и артефактами.

    \item \textbf{Дисциплинированность}: проектирование, кодирование, тестирование, документирование и сопровождение выполняются не произвольно, а в рамках установленных методов, стандартов и соглашений.

    \item \textbf{Измеримость}: состояние проекта и качество результата должны по возможности оцениваться количественно: сроками, стоимостью, числом дефектов, покрытием тестами, производительностью, надежностью, сопровождаемостью и другими метриками.
\end{enumerate}

Программная инженерия охватывает не только собственно написание исходного кода, но и полный жизненный цикл программного обеспечения:

\begin{itemize}
    \item анализ предметной области;
    \item выявление и спецификацию требований;
    \item проектирование архитектуры;
    \item детальное проектирование;
    \item реализацию;
    \item тестирование и верификацию;
    \item внедрение;
    \item эксплуатацию;
    \item сопровождение;
    \item управление конфигурациями;
    \item управление проектом;
    \item обеспечение качества;
    \item вывод программной системы из эксплуатации.
\end{itemize}

Таким образом, программная инженерия занимается не только вопросом \emph{как написать программу}, но и вопросами:

\begin{itemize}
    \item как понять, что именно нужно разработать;
    \item как доказать или проверить, что разработанное соответствует требованиям;
    \item как организовать работу коллектива разработчиков;
    \item как управлять изменениями;
    \item как снизить стоимость сопровождения;
    \item как обеспечить надежность, безопасность и масштабируемость;
    \item как сделать программное обеспечение пригодным для длительного использования.
\end{itemize}

\subsection{Понятие технологии программирования}

\textbf{Технология программирования} --- это совокупность методов, приемов, языков, инструментов, правил и процессов, используемых для создания программ.

В более узком смысле технология программирования описывает практическую сторону разработки:

\begin{itemize}
    \item выбор языка программирования;
    \item стиль программирования;
    \item правила структурирования программ;
    \item использование библиотек и фреймворков;
    \item методы отладки;
    \item методы тестирования;
    \item средства сборки;
    \item средства контроля версий;
    \item инструменты автоматизации разработки.
\end{itemize}

Технология программирования может рассматриваться как часть программной инженерии. Если программная инженерия задает общую инженерную рамку создания программных систем, то технология программирования описывает конкретные способы реализации программ в рамках этой рамки.

\subsection{Сравнение программной инженерии и технологии программирования}

\begin{longtable}{|p{4cm}|p{5.5cm}|p{5.5cm}|}
\hline
\textbf{Критерий} & \textbf{Программная инженерия} & \textbf{Технология программирования} \\
\hline
Объект рассмотрения &
Программное обеспечение как инженерная система на протяжении всего жизненного цикла &
Процесс создания программного кода и связанных с ним артефактов \\
\hline
Масштаб &
Широкий: требования, проектирование, реализация, тестирование, сопровождение, управление, качество &
Более узкий: методы и средства непосредственной разработки программ \\
\hline
Основной вопрос &
Как организовать надежное, управляемое и экономически оправданное создание программных систем? &
Как эффективно реализовать программу? \\
\hline
Типичные темы &
Жизненный цикл ПО, управление проектами, требования, архитектура, качество, сопровождение, стандарты &
Парадигмы программирования, языки, отладка, тестирование, библиотеки, инструменты разработки \\
\hline
Результат &
Программная система, удовлетворяющая требованиям заказчика и ограничениям проекта &
Программный код и связанные с ним компоненты \\
\hline
\end{longtable}

\subsection{Появление программной инженерии как науки}

Программная инженерия возникла как ответ на практические трудности разработки крупных программных систем. По мере роста вычислительной техники программное обеспечение становилось все более сложным, дорогим и критически важным. Методы индивидуального программирования, пригодные для небольших задач, оказались недостаточными для создания больших систем.

Исторически важным понятием стала \textbf{кризис программного обеспечения}. Под ним понимают совокупность проблем, возникших при разработке сложных программных систем:

\begin{itemize}
    \item проекты часто завершались с опозданием;
    \item стоимость разработки превышала первоначальные оценки;
    \item программы содержали большое количество ошибок;
    \item программные системы было трудно сопровождать;
    \item фактическое поведение программ не соответствовало ожиданиям пользователей;
    \item документация была неполной или устаревшей;
    \item изменение требований приводило к неконтролируемому росту сложности;
    \item разработка зависела от отдельных программистов и была плохо воспроизводимой.
\end{itemize}

\subsection{Основные причины возникновения программной инженерии}

\subsubsection{Рост сложности программных систем}

Первые программы были сравнительно небольшими и часто создавались одним человеком. С развитием вычислительной техники появились операционные системы, компиляторы, системы управления базами данных, автоматизированные системы управления, банковские системы, телекоммуникационные системы, встроенное ПО, распределенные и веб-системы.

Такие системы обладают следующими особенностями:

\begin{itemize}
    \item состоят из большого числа компонентов;
    \item разрабатываются коллективами специалистов;
    \item работают в течение многих лет;
    \item взаимодействуют с другими системами;
    \item требуют высокой надежности;
    \item должны изменяться вместе с изменением предметной области.
\end{itemize}

Сложность программных систем имеет несколько аспектов:

\begin{enumerate}
    \item \textbf{Функциональная сложность}: большое количество функций и сценариев использования.
    \item \textbf{Структурная сложность}: большое количество модулей, классов, интерфейсов и зависимостей.
    \item \textbf{Алгоритмическая сложность}: необходимость реализации нетривиальных алгоритмов обработки данных.
    \item \textbf{Организационная сложность}: участие многих разработчиков, аналитиков, тестировщиков и менеджеров.
    \item \textbf{Эволюционная сложность}: постоянное изменение требований и среды эксплуатации.
\end{enumerate}

\subsubsection{Высокая стоимость ошибок}

Ошибки в программном обеспечении могут приводить не только к неудобствам, но и к серьезным экономическим, техническим и социальным последствиям.

Примеры критичных областей:

\begin{itemize}
    \item авиация;
    \item медицина;
    \item атомная энергетика;
    \item банковские системы;
    \item системы управления транспортом;
    \item телекоммуникации;
    \item промышленная автоматизация;
    \item информационная безопасность.
\end{itemize}

В таких областях программное обеспечение должно удовлетворять строгим требованиям к надежности, безопасности, отказоустойчивости и проверяемости.

\subsubsection{Необходимость коллективной разработки}

Крупные программные системы невозможно эффективно разрабатывать без разделения труда. Коллективная разработка требует:

\begin{itemize}
    \item единого понимания требований;
    \item согласованной архитектуры;
    \item правил кодирования;
    \item управления версиями;
    \item документирования;
    \item тестирования;
    \item планирования задач;
    \item управления изменениями;
    \item контроля качества.
\end{itemize}

Без инженерной организации коллективная разработка приводит к хаосу: дублированию работы, несовместимости модулей, потере изменений, росту числа дефектов и невозможности сопровождения.

\subsubsection{Проблема сопровождения программ}

Сопровождение --- это изменение программного обеспечения после передачи в эксплуатацию. Оно может включать:

\begin{itemize}
    \item исправление ошибок;
    \item адаптацию к новой аппаратной или программной среде;
    \item добавление новых функций;
    \item улучшение производительности;
    \item повышение безопасности;
    \item рефакторинг;
    \item обновление документации.
\end{itemize}

Значительная часть затрат на программную систему возникает не во время первоначальной разработки, а во время ее эксплуатации и сопровождения. Поэтому важно проектировать программы так, чтобы они были понятными, модульными, тестируемыми и изменяемыми.

\subsubsection{Необходимость управления качеством}

Качество программного обеспечения нельзя обеспечить только финальным тестированием. Оно должно закладываться на всех этапах жизненного цикла:

\begin{itemize}
    \item при формулировании требований;
    \item при выборе архитектуры;
    \item при проектировании интерфейсов;
    \item при реализации;
    \item при тестировании;
    \item при внедрении;
    \item при сопровождении.
\end{itemize}

Качество ПО включает не только отсутствие ошибок, но и множество характеристик:

\begin{itemize}
    \item функциональная пригодность;
    \item надежность;
    \item удобство использования;
    \item производительность;
    \item безопасность;
    \item совместимость;
    \item сопровождаемость;
    \item переносимость.
\end{itemize}

\subsubsection{Потребность в стандартизации и повторяемости}

Инженерная дисциплина предполагает, что успешные методы разработки должны быть описаны, изучены, воспроизведены и улучшены. Поэтому в программной инженерии большое значение имеют:

\begin{itemize}
    \item стандарты;
    \item методологии;
    \item модели жизненного цикла;
    \item шаблоны проектирования;
    \item практики управления требованиями;
    \item практики тестирования;
    \item модели зрелости процессов;
    \item метрики качества и процесса.
\end{itemize}

\section{Программное изделие и программный продукт}

\subsection{Понятие программного изделия}

\textbf{Программное изделие} --- это завершенный или подлежащий поставке результат разработки программного обеспечения, оформленный как самостоятельная единица, имеющая состав, назначение, документацию, версию и правила эксплуатации.

Программное изделие обычно включает не только исполняемый код, но и связанные с ним компоненты:

\begin{itemize}
    \item исходные тексты;
    \item исполняемые файлы;
    \item конфигурационные файлы;
    \item базы данных или схемы данных;
    \item пользовательскую документацию;
    \item эксплуатационную документацию;
    \item инструкции по установке;
    \item тестовые материалы;
    \item лицензионные соглашения;
    \item средства сопровождения и обновления.
\end{itemize}

Термин \emph{изделие} подчеркивает инженерный и производственный характер результата. Программное изделие может рассматриваться аналогично техническому изделию: оно имеет состав, характеристики, версию, документацию, условия применения и правила приемки.

\subsection{Понятие программного продукта}

\textbf{Программный продукт} --- это программное обеспечение, предназначенное для поставки, распространения, продажи, внедрения или использования некоторым кругом пользователей.

Программный продукт ориентирован не только на техническую готовность, но и на потребительскую ценность. Он должен решать задачу пользователя и быть пригодным для практического применения.

Программный продукт обычно характеризуется следующими признаками:

\begin{itemize}
    \item имеет целевую аудиторию;
    \item удовлетворяет определенным потребностям пользователей;
    \item сопровождается документацией;
    \item имеет механизм установки или доступа;
    \item имеет версионность;
    \item может иметь лицензию;
    \item может иметь службу поддержки;
    \item может развиваться в соответствии с запросами рынка или заказчика.
\end{itemize}

\subsection{Соотношение программного изделия и программного продукта}

Понятия \emph{программное изделие} и \emph{программный продукт} близки, но акцентируют разные стороны программного обеспечения.

\begin{itemize}
    \item \textbf{Программное изделие} --- акцент на инженерной завершенности, составе, документации, контроле версий и поставке.
    \item \textbf{Программный продукт} --- акцент на потребительской ценности, использовании, распространении и удовлетворении потребностей пользователей.
\end{itemize}

Можно сказать, что программный продукт обычно является программным изделием, но не каждое программное изделие обязательно является рыночным продуктом. Например, внутренний программный модуль предприятия может быть изделием в рамках проекта, но не быть продуктом для внешнего рынка.

\subsection{Состав программного продукта}

Типичный программный продукт включает:

\begin{enumerate}
    \item \textbf{Функциональную часть}: реализованные функции, ради которых продукт создается.

    \item \textbf{Интерфейс пользователя}: графический, командный, программный или иной интерфейс взаимодействия.

    \item \textbf{Документацию}: пользовательскую, техническую, эксплуатационную, административную.

    \item \textbf{Средства установки и обновления}: инсталляторы, контейнеры, пакеты, скрипты развертывания.

    \item \textbf{Средства диагностики}: журналы, механизмы мониторинга, отчеты об ошибках.

    \item \textbf{Лицензионные и правовые материалы}: лицензия, условия использования, сведения об авторских правах.

    \item \textbf{Средства сопровождения}: механизмы обратной связи, система поддержки, база знаний, процедуры обновления.
\end{enumerate}

\section{Типы программных приложений}

\subsection{Промышленное программное приложение}

\textbf{Промышленное программное приложение} --- это программное обеспечение, создаваемое для практического использования в реальных условиях эксплуатации, обычно с учетом требований заказчика, рынка, нормативных ограничений, надежности, безопасности, сопровождаемости и экономической эффективности.

Промышленные приложения могут быть:

\begin{itemize}
    \item заказными;
    \item коробочными;
    \item облачными;
    \item встроенными;
    \item корпоративными;
    \item массовыми пользовательскими;
    \item критически важными.
\end{itemize}

Основные признаки промышленного приложения:

\begin{enumerate}
    \item \textbf{Ориентация на пользователя или заказчика}. Приложение создается для решения практической задачи.

    \item \textbf{Требования к надежности}. Программа должна устойчиво работать в условиях реальной эксплуатации.

    \item \textbf{Документированность}. Необходимы пользовательская, техническая и эксплуатационная документация.

    \item \textbf{Поддержка и сопровождение}. После выпуска приложение исправляется, обновляется и развивается.

    \item \textbf{Контроль качества}. Используются тестирование, ревью, статический анализ, CI/CD, мониторинг.

    \item \textbf{Управление версиями и релизами}. Выпуски продукта планируются, маркируются и сопровождаются изменениями.

    \item \textbf{Экономические ограничения}. Разработка связана с бюджетом, сроками и коммерческими рисками.

    \item \textbf{Юридическая определенность}. Обычно существуют договоры, лицензии, SLA, гарантии или ограничения ответственности.
\end{enumerate}

\subsection{Open source программное приложение}

\textbf{Open source программное приложение} --- это программное обеспечение, исходный код которого открыт для просмотра, изучения, изменения и распространения в соответствии с условиями открытой лицензии.

Открытость исходного кода не означает отсутствия авторских прав. Напротив, права регулируются лицензией, например:

\begin{itemize}
    \item MIT License;
    \item Apache License 2.0;
    \item GNU GPL;
    \item LGPL;
    \item BSD License;
    \item Mozilla Public License.
\end{itemize}

Основные признаки Open source приложения:

\begin{enumerate}
    \item \textbf{Доступность исходного кода}. Пользователи и разработчики могут изучать реализацию.

    \item \textbf{Право модификации}. Лицензия обычно разрешает изменять код.

    \item \textbf{Право распространения}. Лицензия определяет условия распространения оригинальной и измененной версии.

    \item \textbf{Сообщество разработчиков}. Развитие часто происходит через внешние вклады, pull request'ы, обсуждения и issue tracker.

    \item \textbf{Публичная история изменений}. Используются открытые репозитории и системы контроля версий.

    \item \textbf{Разная степень формальности процессов}. Некоторые проекты имеют строгую инженерную дисциплину, другие развиваются менее формально.
\end{enumerate}

Open source проект может быть как любительским, так и промышленным. Например, многие открытые проекты используются в критически важных промышленных системах и поддерживаются крупными организациями.

\subsection{Исследовательское программное приложение}

\textbf{Исследовательское программное приложение} --- это программное обеспечение, создаваемое прежде всего для проверки научной гипотезы, демонстрации метода, проведения эксперимента или получения исследовательских результатов.

Главная цель исследовательского ПО --- не всегда промышленная эксплуатация, а получение нового знания.

Примеры исследовательского ПО:

\begin{itemize}
    \item прототип алгоритма;
    \item экспериментальная система машинного обучения;
    \item симулятор физического процесса;
    \item программа для обработки экспериментальных данных;
    \item доказательство концепции;
    \item демонстрационная реализация метода из статьи.
\end{itemize}

Основные признаки исследовательского приложения:

\begin{enumerate}
    \item \textbf{Ориентация на эксперимент}. Программа создается для проверки идеи или гипотезы.

    \item \textbf{Ограниченный круг пользователей}. Часто пользователями являются сами исследователи или небольшая научная группа.

    \item \textbf{Приоритет научной новизны}. Важнее показать возможность метода, чем обеспечить промышленную надежность.

    \item \textbf{Неполная документация}. Документация может ограничиваться статьей, README или комментариями.

    \item \textbf{Ограниченная сопровождаемость}. Код может не развиваться после завершения эксперимента.

    \item \textbf{Экспериментальная архитектура}. Возможны упрощения, допущения и отсутствие обработки всех исключительных ситуаций.
\end{enumerate}

\section{Сравнение промышленного, Open source и исследовательского приложения}

\begin{longtable}{|p{3.6cm}|p{4cm}|p{4cm}|p{4cm}|}
\hline
\textbf{Критерий} &
\textbf{Промышленное приложение} &
\textbf{Open source приложение} &
\textbf{Исследовательское приложение} \\
\hline
Основная цель &
Решение практической задачи пользователя или заказчика &
Создание открытого ПО, доступного для изучения, изменения и распространения &
Проверка гипотезы, метода или научной идеи \\
\hline
Ключевой результат &
Стабильный продукт для эксплуатации &
Открытый код и работающая система, развиваемая авторами или сообществом &
Прототип, экспериментальная реализация или инструмент исследования \\
\hline
Пользователи &
Клиенты, организации, массовые пользователи, внутренние подразделения &
Сообщество, разработчики, компании, конечные пользователи &
Исследователи, научные группы, рецензенты, участники эксперимента \\
\hline
Требования к надежности &
Обычно высокие, особенно в критичных областях &
Зависят от зрелости проекта и области применения &
Часто ниже, если надежность не является предметом исследования \\
\hline
Документация &
Обычно обязательна: пользовательская, техническая, эксплуатационная &
Может быть хорошей, но зависит от сообщества и зрелости проекта &
Часто минимальна: статья, README, комментарии \\
\hline
Сопровождение &
Плановое, договорное или коммерчески обусловленное &
Может выполняться сообществом, мейнтейнерами или компаниями &
Может отсутствовать после завершения исследования \\
\hline
Качество кода &
Обычно контролируется внутренними стандартами и процессами &
Зависит от правил проекта, ревью и активности сообщества &
Может уступать промышленному уровню из-за экспериментального характера \\
\hline
Лицензирование &
Коммерческое, проприетарное, корпоративное или смешанное &
Открытая лицензия &
Может быть открытым, закрытым или неформально распространяемым \\
\hline
Процесс разработки &
Формализованный: требования, планирование, тестирование, релизы &
Публичный или частично публичный; часто через репозиторий и issue tracker &
Гибкий, экспериментальный, ориентированный на быстрые изменения \\
\hline
Критерий успешности &
Удовлетворение требований, прибыльность, надежность, внедрение &
Полезность, распространенность, активность сообщества, качество вкладов &
Научный результат, публикация, подтверждение или опровержение гипотезы \\
\hline
\end{longtable}

\section{Жизненный цикл программного обеспечения}

Для связного понимания различий между видами программных приложений полезно рассмотреть понятие жизненного цикла ПО.

\textbf{Жизненный цикл программного обеспечения} --- это совокупность стадий, через которые проходит программная система от возникновения идеи до прекращения эксплуатации.

Типовые стадии жизненного цикла:

\begin{enumerate}
    \item \textbf{Инициация}. Определяется проблема, цель разработки и заинтересованные стороны.

    \item \textbf{Анализ требований}. Выявляются функциональные и нефункциональные требования.

    \item \textbf{Проектирование}. Определяется архитектура, структура компонентов, интерфейсы и данные.

    \item \textbf{Реализация}. Создается исходный код и связанные с ним артефакты.

    \item \textbf{Тестирование}. Проверяется соответствие реализации требованиям.

    \item \textbf{Внедрение}. Система устанавливается или публикуется для пользователей.

    \item \textbf{Эксплуатация}. Система используется по назначению.

    \item \textbf{Сопровождение}. Вносятся исправления, улучшения и адаптационные изменения.

    \item \textbf{Вывод из эксплуатации}. Система заменяется, архивируется или прекращает использоваться.
\end{enumerate}

Для промышленного ПО все стадии обычно явно выражены. Для Open source ПО жизненный цикл может быть публичным и распределенным. Для исследовательского ПО некоторые стадии могут быть сокращены или отсутствовать, если они не нужны для достижения научной цели.

\section{Качество программного обеспечения}

\subsection{Понятие качества ПО}

\textbf{Качество программного обеспечения} --- это степень, в которой программная система удовлетворяет установленным и предполагаемым потребностям пользователей, заказчиков и других заинтересованных сторон.

Качество не сводится только к корректности. Программа может правильно выполнять вычисления, но быть неудобной, медленной, небезопасной или практически несопровождаемой.

\subsection{Основные характеристики качества}

К основным характеристикам качества ПО относятся:

\begin{enumerate}
    \item \textbf{Функциональная пригодность}: способность выполнять требуемые функции.

    \item \textbf{Надежность}: способность сохранять работоспособность в течение времени и при наличии сбоев.

    \item \textbf{Производительность}: эффективность использования времени, памяти, сети и других ресурсов.

    \item \textbf{Удобство использования}: простота освоения и применения пользователем.

    \item \textbf{Безопасность}: защита данных и функций от несанкционированного доступа и повреждения.

    \item \textbf{Совместимость}: способность взаимодействовать с другими системами.

    \item \textbf{Сопровождаемость}: удобство анализа, изменения, тестирования и развития системы.

    \item \textbf{Переносимость}: возможность переноса в другую программную или аппаратную среду.
\end{enumerate}

\section{Программная инженерия как инженерная дисциплина}

Программная инженерия отличается от индивидуального программирования тем, что рассматривает создание ПО как управляемый инженерный процесс. Это означает:

\begin{itemize}
    \item наличие требований;
    \item планирование работ;
    \item оценку ресурсов;
    \item управление рисками;
    \item проектирование архитектуры;
    \item применение стандартов;
    \item контроль качества;
    \item сопровождение результата;
    \item учет стоимости владения системой.
\end{itemize}

В инженерном подходе важны не только творческие решения, но и воспроизводимость. Хороший процесс разработки должен позволять разным специалистам получать предсказуемый результат при сходных условиях.

\section{Краткие выводы}

\begin{enumerate}
    \item Программная инженерия возникла как ответ на рост сложности программных систем и кризис традиционных методов индивидуального программирования.

    \item Программная инженерия шире технологии программирования: она охватывает весь жизненный цикл ПО, управление проектами, требования, качество, сопровождение и экономические аспекты.

    \item Технология программирования описывает методы и средства непосредственного создания программ: языки, инструменты, стили, библиотеки, отладку и тестирование.

    \item Программное изделие --- инженерно оформленный результат разработки, имеющий состав, версию, документацию и назначение.

    \item Программный продукт --- программное обеспечение, ориентированное на использование пользователями и обладающее потребительской ценностью.

    \item Промышленное приложение ориентировано на практическую эксплуатацию, надежность, сопровождение и выполнение требований заказчика или рынка.

    \item Open source приложение отличается открытостью исходного кода и правом его изучения, изменения и распространения в рамках лицензии.

    \item Исследовательское приложение создается прежде всего для проверки научной идеи, метода или гипотезы; его промышленное качество не всегда является основной целью.
\end{enumerate}

\end{document}
